\chapter{Носитель}
\label{ch:carrier-object}

\section{Место для состояния}
\label{sec:carrier-object}

В гл.~\ref{ch:descent} зафиксирован вопрос, а не модель:
наблюдаемая физика упирается в планковский предел,
и любой более глубокий закон обязан её воспроизвести.
Чтобы сказать этот закон, нужно сначала назвать, \emph{где} живёт состояние.

Настоящая глава вводит только этот объект.
Счётное множество узлов, дискретное время, элементарные шаг и такт,
окрестность одного такта и состояние, приписанное узлу.
Какой граф связывает узлы, сколько у узла соседей
и чему равна безразмерная константа между тактовой скоростью и наблюдаемой скоростью света,
здесь не выбирается.
Это теоремы гл.~\ref{ch:carrier}: они становятся возможны только после того,
как на языке этой главы будут записаны абсолютные условия.

Сами символы $\Lambda$, $\hl$, $\htick$ пока средство речи, а не отдельный постулат
«решётка существует''.
Что континуум на микроуровне невозможен и что шаг нельзя устремить к нулю,
доказывается позже (гл.~\ref{ch:foundations}), когда есть и условия, и геометрия.

\section{Узлы, шаг и такт}
\label{sec:sites-step-tick}

\begin{definition}[Носитель как множество узлов]
\label{def:carrier-sites}
Носитель микроуровня~$M$ --- счётное множество узлов~$\Lambda$.
Время дискретно: $t\in\mathbb{Z}$.
Граф, рёбра которого соединяют узлы, в это определение не входит.
\end{definition}

\begin{definition}[Элементарные шаг и такт]
\label{def:step-tick}
$\hl$ --- элементарный шаг: единственная длина ребра, на которое за один такт может уйти сигнал.
$\htick$ --- один такт эволюции.
Тактовая скорость сигнала на~$M$
\[
  c_0 = \frac{\hl}{\htick}.
\]
Макроскопическая скорость света~$c$ с~$c_0$ здесь не отождествляется.
Безразмерный множитель в связи $c=\kgeom c_0$ появится вместе с геометрией
однотактового тела (гл.~\ref{ch:carrier}).
Отождествление $\hl$ с планковской длиной --- в гл.~\ref{ch:si-sm}.
\end{definition}

\begin{definition}[Окрестность одного такта]
\label{def:neighborhood}
Окрестность узла $x\in\Lambda$ --- множество узлов, достижимых из~$x$
ровно за один такт по ребру длины~$\hl$:
\[
  N(x)\subset\Lambda.
\]
Число $|N(x)|$ и тип решётки определение не задаёт.
Световая формула окрестности --- содержание условия каузальности;
единственный граф в трёх измерениях --- теорема гл.~\ref{ch:carrier}.
\end{definition}

\section{Состояние на узле}
\label{sec:state-on-carrier}

\begin{definition}[Поле и норма]
\label{def:field-a3}
\textbf{Узел регистра.}
Элементарная ячейка~$\PlanckCell$ --- \emph{один} узел~$\Lambda$
(масштаб длины~$\hl$), а не точка континуума и не макроскопическая частица.

\textbf{Спинор на ячейке.}
На каждом~$x$ задан \emph{спинор}
\[
  z(x,t)=(z_1,z_2)^\top\in\mathbb{C}^2,
\]
то есть столбец из двух комплексных чисел, а не одно число $z\in\mathbb{C}$.
Две компоненты нужны, чтобы на одном узле были возможны заряд ($U(1)$-фаза),
локальный поворот ($SU(2)$) и спин~$1/2$:
оборот на $2\pi$ меняет знак состояния, оборот на $4\pi$ возвращает его.
Модуль $|z_1|^2+|z_2|^2$ --- плотность поля на ячейке;
фазы $z_k/|z_k|$ --- степени свободы, которые перемешиваются при эволюции.

\textbf{Глобальное состояние.}
Конфигурация всей решётки в такт~$t$:
\[
  \Psi^t=\{z(x,t)\}_{x\in\Lambda}.
\]

\textbf{Норма.}
Суммарная информационная норма (аналог $\langle\psi|\psi\rangle$ в квантовой механике):
\[
  N(t):=\sum_{x\in\Lambda} z^\dagger(x,t)\,z(x,t)=\sum_x\bigl(|z_1|^2+|z_2|^2\bigr).
\]
Сохраняется ли она при $t\to t+1$, решают условия следующей главы, а не это определение.
\end{definition}

\section{Что этой главой не решено}

Не выбрана решётка и не названо число соседей.
Не записан закон~$g$.
Не доказано, что иной носитель невозможен:
сначала условия на допустимый локальный закон
(гл.~\ref{ch:axiom-rationale}, \ref{ch:axioms}),
затем геометрия, которая из них следует (гл.~\ref{ch:carrier}),
и уже после этого --- теорема дискретности (гл.~\ref{ch:foundations}).
